\bibitem{Maldacena}
J.~M.~Maldacena, \emph{Wilson loops in large $N$ field theories},
Physical Review Letters \textbf{80} (1998), 4859--4862,
\href{https://arxiv.org/abs/hep-th/9803002}{arXiv:hep-th/9803002}.

\bibitem{Zarembo}
K.~Zarembo, \emph{Supersymmetric Wilson loops},
Nuclear Physics B \textbf{643} (2002), 157--171,
\href{https://arxiv.org/abs/hep-th/0205160}{arXiv:hep-th/0205160}.

\bibitem{DFO}
O.~DeWolfe, D.~Z.~Freedman and H.~Ooguri,
\emph{Holography and defect conformal field theories},
Physical Review D \textbf{66} (2002), 025009,
\href{https://arxiv.org/abs/hep-th/0111135}{arXiv:hep-th/0111135}.

\bibitem{GW}
D.~Gaiotto and E.~Witten,
\emph{Supersymmetric boundary conditions in $\mathcal N=4$ super
Yang--Mills theory}, Journal of Statistical Physics \textbf{135}
(2009), 789--855,
\href{https://arxiv.org/abs/0804.2902}{arXiv:0804.2902}.

\bibitem{Baker}
E.~B.~Baker~III, \emph{Supersymmetric open Wilson lines},
Journal of High Energy Physics \textbf{03} (2011), 140,
\href{https://arxiv.org/abs/1102.4948}{arXiv:1102.4948}.

\bibitem{Lawler}
G.~F.~Lawler, \emph{Schramm--Loewner evolution (SLE)}, lecture notes,
especially Lecture~2, Proposition~2.17 and equation~(2.8),
\href{https://math.uchicago.edu/~lawler/utah.pdf}{author's PDF},
accessed 19 September 2026.

\bibitem{BauerBernard}
M.~Bauer and D.~Bernard,
\emph{Conformal field theories of stochastic Loewner evolutions},
Communications in Mathematical Physics \textbf{239} (2003), 493--521,
\href{https://arxiv.org/abs/hep-th/0210015}{arXiv:hep-th/0210015}.

\bibitem{OS}
K.~Osterwalder and R.~Schrader,
\emph{Axioms for Euclidean Green's functions},
Communications in Mathematical Physics \textbf{31} (1973), 83--112,
\href{https://doi.org/10.1007/BF01645738}{doi:10.1007/BF01645738}.

\bibitem{DP}
H.~Dorn and V.~D.~Pershin,
\emph{Concavity of the $Q\bar Q$ potential in $\mathcal N=4$ super
Yang--Mills gauge theory and AdS/CFT duality},
Physics Letters B \textbf{461} (1999), 338--344,
\href{https://arxiv.org/abs/hep-th/9906073v3}{arXiv:hep-th/9906073v3}.

\bibitem{Suzuki}
M.~Suzuki, \emph{Weil's quadratic form via the screw function},
preprint (2026), version~2,
\href{https://arxiv.org/abs/2606.09096v2}{arXiv:2606.09096v2}.
Used for the classical compact-support Weil criterion and conventions;
its proposed limiting spectral construction is not assumed here.

\bibitem{ParentWilson}
Riemann Hypothesis research program, \emph{Deformation flows for the
shifted Weil family: the transfer as an all-pass filter}, working
draft~0.8, 17 September 2026, and endpoint continuation notes,
18--19 September 2026. Repository investigation
\nolinkurl{papers/susy-positivity/investigations/wilson-lines/}.
Internal research, with model-assisted contributions disclosed in
the source documents.

\bibitem{ParentAngular}
OpenAI GPT-6 (Codex), for Edward Baker,
\emph{Angular smearing, the finite difference form, and a Robin model
of the shift}, research note, 19 September 2026, in the Wilson-lines
investigation, \nolinkurl{notes/ANGULAR_SMEARING_AND_ROBIN_MODEL_20260919.md}.

\bibitem{ParentLoewner}
Riemann Hypothesis research program,
\emph{The Markov part of the shifted Weil transfer}, working
draft~0.3, 18 September 2026. Repository investigation
\nolinkurl{papers/susy-positivity/investigations/loewner/}.
Internal research; its direct realization proposal was closed within
the tested scope.

\bibitem{ParentFractional}
Riemann Hypothesis research program,
\emph{A scattering matrix without a space}, working draft~0.1,
18 September 2026. Repository investigation
\nolinkurl{papers/susy-positivity/investigations/fractional-dimension/}.
Internal research; the continued lattice count and its limitations
are discussed in its Section~6.


\bibitem{SharedBackground}
Riemann Hypothesis research program,
\emph{Shared background for the supersymmetric positivity program},
repository source
\nolinkurl{papers/susy-positivity/background_section.tex},
state consulted 19 September 2026. In particular, the boxed
transfer and defect equations beside the energy derivative
and the sufficient depth--shift implication.

\bibitem{ShiftFlowNote}
OpenAI GPT-6 (Codex), for Edward Baker,
\emph{Prescribed shift evolution, cumulative storage, and a
depth--shift continuation}, research note, 19 September 2026,
\nolinkurl{notes/SHIFT_FLOW_CUMULATIVE_STORAGE_AND_CRITICAL_PATH_20260919.md}
in the Wilson--Loewner investigation.

\bibitem{AdaptivePath}
Riemann Hypothesis research program,
\emph{Depth extension with an adaptive shift}, working research,
10 September 2026,
\nolinkurl{papers/shifted-zeta/storage-depth/archive/background/1-critical_path_research.md}.
Internal, model-assisted research; independent review remains open.

\bibitem{CumulativePath}
Riemann Hypothesis research program,
\emph{A depth--shift path controlled by cumulative storage},
working research, 10 September 2026,
\nolinkurl{papers/shifted-zeta/storage-depth/archive/background/2-cumulative_storage_path.md}.
Internal, model-assisted research; the continuation bounds are
conditional.

\bibitem{StorageDepth}
E.~Baker, \emph{Residual-controlled depth extension of finite-horizon
Weil positivity}, working draft~0.3, with model assistance disclosed
there; see Appendix~A,
\nolinkurl{papers/shifted-zeta/storage-depth/manuscript/storage_depth.tex}.
The recorded enclosures use that paper's working normalization;
independent normalization verification and review remain open.

\bibitem{Suzuki2012}
M.~Suzuki, \emph{A canonical system of differential equations arising
from the Riemann zeta-function}, RIMS Kokyuroku Bessatsu
\textbf{B34} (2012), 397--435, especially Section~1.4,
\href{https://www.kurims.kyoto-u.ac.jp/~kenkyubu/bessatsu/open/B34/pdf/B34_023.pdf}{publisher's PDF}.


\bibitem{Pestun}
V.~Pestun, \emph{Localization of gauge theory on a four-sphere and
supersymmetric Wilson loops},
\href{https://arxiv.org/abs/0712.2824}{arXiv:0712.2824},
especially Section~4, equations~(4.23)--(4.25).

\bibitem{WangDefects}
Y.~Wang, \emph{Taming defects in $\mathcal N=4$ super-Yang--Mills},
Journal of High Energy Physics \textbf{08} (2020), 021,
\href{https://arxiv.org/abs/2003.11016}{arXiv:2003.11016},
especially Sections~4.4 and 5.1, equation~(5.5).

\bibitem{DLMFGamma}
NIST Digital Library of Mathematical Functions,
Chapter~5, \emph{Gamma function},
\href{https://dlmf.nist.gov/5.5}{Section~5.5} (recurrences) and
\href{https://dlmf.nist.gov/5.11}{Section~5.11} (asymptotics),
accessed 20 September 2026.

\bibitem{DPYHiggs}
M.~Dedushenko, S.~S.~Pufu and R.~Yacoby,
\emph{A one-dimensional theory for Higgs branch operators},
\href{https://arxiv.org/abs/1610.00740v2}{arXiv:1610.00740v2},
especially Sections~3.1 and 6.1.

\bibitem{DedushenkoGluing}
M.~Dedushenko, \emph{Gluing II: Boundary localization and gluing formulas},
\href{https://arxiv.org/abs/1807.04278v3}{arXiv:1807.04278v3},
especially Sections~4.2 and 4.3.2.

\bibitem{EMAAnchor}
OpenAI GPT-6 (Codex), for Edward Baker,
\emph{The arithmetic EMA tower gives an all-input anchor},
research note, 20 September 2026,
\nolinkurl{notes/EMA_TOWER_ORIGINAL_TRANSFER_ANCHOR_20260920.md}
in the neighboring critical-path investigation.
Computer-assisted fixed-window proof with rational interval records;
independent specialist review remains outstanding.

\bibitem{CumulativeAppend}
OpenAI GPT-6 (Codex), for Edward Baker,
\emph{Cumulative append: scalar-floor obstruction and finite diagnostics},
research note, 20 September 2026,
\nolinkurl{notes/CUMULATIVE_APPEND_SCALAR_OBSTRUCTION_20260920.md}
in the neighboring critical-path investigation.
This records the scalar estimate and finite diagnostics; the later
operator-coupling calculation is cited separately below.

\bibitem{CumulativeCoupling}
OpenAI GPT-6 (Codex), for Edward Baker,
\emph{An all-input cumulative append bound from forward/backward energy},
research note, 20 September 2026,
\nolinkurl{notes/CUMULATIVE_OPERATOR_COUPLING_CERTIFICATE_20260920.md}
in the neighboring critical-path investigation.
Internal computer-assisted proof; independent specialist review outstanding.

\bibitem{GrowingTraceNote}
OpenAI GPT-6 (Codex), for Edward Baker,
\emph{Genuine Loewner growth in fixed Yang--Mills theory: the first two
quantum evolution equations}, research note, 21 September 2026,
\nolinkurl{notes/GENUINE_LOEWNER_TRACE_YM_HIERARCHY_20260921.md}.
Internal derivation and finite transport checks under stated quantum assumptions.

\bibitem{LuscherFlow}
M.~L\"uscher, \emph{Properties and uses of the Wilson flow in lattice QCD},
\href{https://arxiv.org/abs/1006.4518}{arXiv:1006.4518}.

\bibitem{LuscherWeiszFlow}
M.~L\"uscher and P.~Weisz, \emph{Perturbative analysis of the gradient
flow in non-abelian gauge theories},
\href{https://arxiv.org/abs/1101.0963}{arXiv:1101.0963}.

\bibitem{LindTran}
J.~Lind and H.~Tran, \emph{Regularity of Loewner curves},
\href{https://arxiv.org/abs/1411.2164}{arXiv:1411.2164}.

\bibitem{MakeenkoLoops}
Y.~Makeenko, \emph{Topics in cusped/lightcone Wilson loops},
\href{https://arxiv.org/abs/0810.2183}{arXiv:0810.2183},
especially Lecture~3, equations~(3.1)--(3.4) and Section~3.6.

\bibitem{ShenSmithZhuLines}
H.~Shen, S.~A.~Smith and R.~Zhu,
\emph{Makeenko--Migdal equations for lattice Yang--Mills--Higgs},
\href{https://arxiv.org/abs/2512.00570}{arXiv:2512.00570}.
